\section{Introduction}

Misleading data visualizations are a recurring issue in financial
reporting. The U.S. Securities and Exchange Commission (SEC) reviews
thousands of 10-K and 10-Q filings annually, where charts are frequently
used to present financial performance. While these visualizations are
intended to improve readability, certain design choices such as
truncated axes, uneven scaling, or three-dimensional effects can
unintentionally, or sometimes intentionally, distort how information is
perceived. Prior work shows that even simple changes like truncating the
y-axis can significantly exaggerate differences and affect user judgment
\cite{pandey2015deceptive}. To reduce these risks, regulations such as
Regulation~G and Item~10(e) of Regulation~S-K provide guidance on how
non-GAAP financial measures should be presented. However, in practice,
most enforcement focuses on textual disclosures rather than visual
representations, and no automated tools currently exist for chart-level
compliance checking. As a result, identifying misleading charts still
relies heavily on manual review, which is not only time-consuming but
also difficult to scale across thousands of filings.

Recent progress in Vision-Language Models (VLMs) has made it possible to
analyze both visual and textual information jointly. Models evaluated in
ChartQA \cite{masry2022chartqa} and approaches like DePlot
\cite{liu2023deplot} show promising results in extracting and reasoning
over chart data. The Misviz benchmark \cite{tonglet2025misviz} provides
the first dataset specifically designed for misleading visualization
detection, covering 12 distinct misleader types. However, most existing
evaluations focus only on binary classification, asking whether a chart
is misleading or not, and per-type detection ability remains largely
unexplored. A separate line of work on tool-augmented language models,
including Toolformer \cite{schick2023toolformer} and ReAct
\cite{yao2023react}, demonstrates that combining language models with
external tools can improve reasoning performance in text-based tasks.
Whether this benefit extends to visual reasoning tasks such as chart
analysis, where additional inputs like OCR outputs may introduce noise
rather than helpful signals, remains an open question.

In this project, we study how well VLMs can detect misleading financial
charts and how different system designs affect their performance. We
focus on three research questions: (RQ1) how accurately VLMs can
identify each of the 12 misleader types defined by \citet{tonglet2025misviz},
and whether model scale matters; (RQ2) whether tool augmentation
improves detection performance; and (RQ3) whether findings from
benchmark datasets generalize to real-world SEC filings.

This paper makes four contributions. First, we present a systematic
comparison of eight pipeline variants across three VLMs on the Misviz
benchmark \cite{tonglet2025misviz}, evaluating both binary and per-type
detection performance across all 12 misleader categories. Second, we
find that injecting tool outputs such as OCR results and rule engine
verdicts directly into VLM prompts consistently degrades performance,
while using the same tools as post-processing vetoes improves results.
This finding challenges common assumptions from tool-augmented LLM work
\cite{yao2023react, schick2023toolformer} and suggests that visual
reasoning tasks require a different integration strategy. Third, we
propose a ViT-based selective veto mechanism \cite{dosovitskiy2021vit}
trained on the Misviz synthetic dataset, which filters false positives
on five high-precision misleader types without additional VLM calls.
Fourth, we validate our best pipeline on real SEC 10-K filings across
ten public companies, comparing findings against actual SEC comment
letters as external ground truth.

